\section{The different and normal upper orbits}

Let $K$ be a complete discretely valued field with perfect residue field,
let $L/K$ be a finite Galois extension, and let $M$ be any intermediate
field. All valuation rings and residue fields below are those of these
actual extensions. Write
\[
 A=\mathcal O_K,\quad B=\mathcal O_M,\quad C=\mathcal O_L,\qquad
 G=\operatorname{Gal}(L/K),\quad H=\operatorname{Gal}(L/M).
\]
The group $H$ is regarded as its actual subgroup of $G$. Put
\[
 n=[M:K],\quad e=e(M/K),\quad f=[k_M:k_K],\quad
 q=e(L/M)=|H_0|,\quad d=d(M/K).
\]
Thus $n=ef$, $|G_0|=eq$, and $e,f,q$ are positive integers. No
normality of $M/K$ or total ramification of $L/K$ is assumed.
The different convention is
$d(M/K)=v_M(\mathfrak D_{B/A})$.

\begin{lemma}[Transitivity of the different ideal]
\label{fw:different-ideal-transitivity}
\source{serre-1979-local-fields:page-0058}
\dcref{Serre, Chapter III, Section 4, Proposition 8}
\uses{fw:different}
\lean{differentIdeal_eq_differentIdeal_mul_differentIdeal}
\mathlibok
The actual different ideals satisfy
\[
 \mathfrak D_{C/A}=\mathfrak D_{C/B}(\mathfrak D_{B/A}C).
\]
\end{lemma}
\begin{proof}
Write $D_{C/B}^{-1}$ for the codifferent, and similarly for the other
two extensions. For every fractional $C$-ideal $J$, trace transitivity
and the defining trace-dual condition give
\[
\begin{aligned}
 J\subseteq D_{C/B}^{-1}
 &\Longleftrightarrow \operatorname{Tr}_{L/M}(J)\subseteq B\\
 &\Longleftrightarrow
 D_{B/A}^{-1}\operatorname{Tr}_{L/M}(J)\subseteq D_{B/A}^{-1}\\
 &\Longleftrightarrow
 \operatorname{Tr}_{M/K}
   (D_{B/A}^{-1}\operatorname{Tr}_{L/M}(J))\subseteq A\\
 &\Longleftrightarrow
 \operatorname{Tr}_{L/K}(D_{B/A}^{-1}J)\subseteq A\\
 &\Longleftrightarrow D_{B/A}^{-1}J\subseteq D_{C/A}^{-1}\\
 &\Longleftrightarrow J\subseteq D_{B/A}D_{C/A}^{-1}.
\end{aligned}
\]
In the second and last equivalences the nonzero fractional ideal
$D_{B/A}$ is invertible. In the third and fifth equivalences, each
argument is an ideal over the appropriate integer ring, so testing its
trace also tests the trace after multiplication by every integer.
These equivalences for all $J$ identify the two fractional ideals
$D_{C/B}^{-1}=D_{B/A}D_{C/A}^{-1}$. Inverting gives the assertion.
This is the trace-dual proof of \cite{Serre1979}.
\end{proof}

\begin{lemma}[Orders of the different in a tower]
\label{fw:different-order-transitivity}
\uses{fw:different-ideal-transitivity,fw:normalized-valuation-scaling,fw:ideal-order-generator}
\lean{FontaineWintenberger.maximalIdeal_emultiplicity_mul, FontaineWintenberger.maximalIdeal_emultiplicity_ne_top, FontaineWintenberger.maximalIdeal_emultiplicity_map, FontaineWintenberger.differentIdeal_order_tower, FontaineWintenberger.differentIdeal_order_toNat_tower}
\leanok
With the preceding actual valuations,
\[
 d(L/K)=d(L/M)+q\,d(M/K).
\]
\end{lemma}
\begin{proof}
Every nonzero ideal of $B$ or $C$ is a power of its maximal ideal.
In particular, choose a generator $b$ of $\mathfrak D_{B/A}$.
Its valuation is $v_M(b)=d(M/K)$. Its image in $C$ has valuation
$v_L(b)=qv_M(b)$ by normalized valuation scaling. Therefore
the extended ideal $\mathfrak D_{B/A}C$ has order $q\,d(M/K)$.
Orders of products of nonzero principal ideals add. Apply this to the
actual different-ideal identity of the preceding lemma.
\end{proof}

\begin{lemma}[A lower-group formula for a subextension different]
\label{fw:subextension-different-lower-sum}
\source{serre-1979-local-fields:page-0069,serre-1979-local-fields:page-0071}
\dcref{Serre, Chapter IV, Section 1, Proposition 2 and Proposition 4}
\uses{fw:different-order-transitivity,fw:hilbert-different,fw:ramification-subgroups}
\lean{FontaineWintenberger.differentIdeal_order_eq_sum_lower_card_of_cutoff, FontaineWintenberger.differentIdeal_order_toReal_eq_sum_lower_card_of_cutoff, FontaineWintenberger.subextension_different_order_mul_inertia_eq_lower_sum, FontaineWintenberger.intermediate_spectral_integer_isGaloisGroup}
\leanok
For every integer $N\geq0$ such that $G_i=1$ for all $i\geq N$,
\[
 qd=\sum_{i=0}^{N-1}(|G_i|-|G_i\cap H|).
\]
\end{lemma}
\begin{proof}
Lower numbering is compatible with subgroups, so $H_i=G_i\cap H$.
The same cutoff $N$ therefore makes both lower filtrations trivial.
Hilbert's different formula in the two Galois extensions $L/K$ and
$L/M$ gives
\[
 d(L/K)=\sum_{i=0}^{N-1}(|G_i|-1),\qquad
 d(L/M)=\sum_{i=0}^{N-1}(|G_i\cap H|-1).
\]
Subtract these two finite sums and use different-order transitivity.
Each summand difference is nonnegative, since $G_i\cap H\subseteq G_i$.
Here $L/M$ is Galois because $L/K$ is Galois; $H$ need not be normal
in $G$.
\end{proof}

\begin{theorem}[The different and the inertia-relative upper index]
\label{fw:subextension-different-inertia-integral}
\source{serre-1979-local-fields:page-0058,serre-1979-local-fields:page-0071,serre-1979-local-fields:page-0080,serre-1979-local-fields:page-0081}
\dcref{Serre, different transitivity, Hilbert's formula, and the Herbrand change of parameter}
\uses{fw:subextension-different-lower-sum,fw:local-inertia-size,fw:herbrand-integer-sum,fw:herbrand-subgroup-integral}
\lean{FontaineWintenberger.subgroup_upperRelIndex_deficit_integral, FontaineWintenberger.subextension_different_order_eq_upperRelIndex_integral, FontaineWintenberger.card_inertia_div_card_subgroup_inertia, FontaineWintenberger.subextension_different_order_eq_ramificationIdx_integral}
\leanok
For $u\geq0$, put
\[
 b_H(u)=[G_0:H_0G^u].
\]
The nonnegative step function $e-b_H(u)$ has bounded support, and
\[
 d=e-1+\int_0^\infty(e-b_H(u))\,du.
\]
The integral may equivalently be taken up to any sufficiently large
finite cutoff. In particular no normality of $H$ is needed.
\end{theorem}
\begin{proof}
For $u\geq0$, the normal subgroup $G^u$ is contained in $G_0$.
Consequently $H_0\subseteq H_0G^u\subseteq G_0$, so
$1\leq b_H(u)\leq[G_0:H_0]=e$. If $G^u=1$, this index is $e$;
eventual triviality therefore gives bounded support.

Choose an integer $N\geq1$ for which both lower filtrations are
trivial at every index at least $N$, and put
$U=\varphi_G(N)$. The Herbrand integral for the actual subgroup $H$
gives
\[
 \int_0^U b_H(u)\,du
 =\varphi_H(\psi_G(U))=\varphi_H(N).
\]
Here $\psi_G(\varphi_G(N))=N$, and the integrand is integrable by
that same subgroup integral theorem. The finite-sum expressions at
the integer $N$ are
\[
 \varphi_G(N)=\frac1{|G_0|}\sum_{i=1}^N|G_i|,\qquad
 \varphi_H(N)=\frac1q\sum_{i=1}^N|G_i\cap H|.
\]
Since $e=|G_0|/q$, subtracting these equalities and integrating the
constant function $e$ yields
\[
 \int_0^U(e-b_H(u))\,du
 = e\varphi_G(N)-\varphi_H(N)
 =\frac1q\sum_{i=1}^N(|G_i|-|G_i\cap H|).
\]
The summand at $N$ is zero. The lower-group different formula says
that the sum from zero to $N-1$ is $qd$; its zeroth summand is
$|G_0|-|H_0|=q(e-1)$. The last display is therefore $d-e+1$.
The integrand vanishes after $U$, so enlarging the cutoff does not
alter this integral. This proves the stated formula using the
already established subgroup Herbrand integral, rather than requiring
a separate change-of-variables theorem.
\end{proof}

\begin{theorem}[The different as an upper-orbit integral]
\label{fw:subextension-different-orbit-integral}
\source{serre-1979-local-fields:page-0058,serre-1979-local-fields:page-0071,serre-1979-local-fields:page-0080,serre-1979-local-fields:page-0081}
\dcref{Serre, different transitivity, Hilbert's formula, and the upper change of parameter}
\uses{fw:subextension-different-inertia-integral,fw:inertia-residue-index,fw:finite-upper-groups}
\lean{FontaineWintenberger.upperOrbitIndex_le_fixingIndex, FontaineWintenberger.upperRamificationGroup_eq_bot_of_herbrand_cutoff, FontaineWintenberger.upperOrbitDeficit_eq_zero_of_herbrand_cutoff, FontaineWintenberger.upperOrbitIndex_monotoneOn, FontaineWintenberger.intervalIntegrable_upperOrbitIndex, FontaineWintenberger.upperOrbitIndex_eq_zeroIndex_mul_upperRelIndex, FontaineWintenberger.fixingIndex_eq_zeroIndex_mul_inertiaRatio, FontaineWintenberger.subextension_different_order_eq_upperOrbitIndex_integral, FontaineWintenberger.intermediate_spectral_different_order_eq_upperOrbitIndex_integral}
\leanok
For $u\geq0$, let
\[
 a_H(u)=[G:HG^u].
\]
This is the number of $G^u$-orbits on the $n$ embeddings of $M$ in
$L$ over $K$. The nonnegative step function $n-a_H(u)$ has bounded
support, and
\[
 f d=n-f+\int_0^\infty(n-a_H(u))\,du.
\]
In particular, if $M/K$ is totally ramified, then
\[
 d=n-1+\int_0^\infty(n-a_H(u))\,du.
\]
\end{theorem}
\begin{proof}
The embedding set identifies with $G/H$. For a normal subgroup
$U\subseteq G$, the stabilizer of $gH$ under $U$ is
$U\cap gHg^{-1}=g(U\cap H)g^{-1}$. Thus every $U$-orbit has size
$|U|/|U\cap H|$. Its orbits correspond exactly to the cosets of
$HU$ in $G$: two cosets $gH,g'H$ are in one orbit exactly when
$g'\in UgH=gHU$. Consequently their number is
\[
 [G:HU]=n\frac{|U\cap H|}{|U|}.
\]
Apply this to the actual normal subgroup $U=G^u$.

We compare this absolute orbit index with the inertia-relative index
$b_H(u)$ of the preceding theorem. The index formula gives
\[
 [G:HG_0]=\frac{[G:H]}{[G_0:H_0]}=\frac ne=f.
\]
For $U\subseteq G_0$, the intersection $G_0\cap HU$ is $H_0U$:
in an expression $hu\in G_0$, the conditions $u\in G_0$ and
$hu\in G_0$ force $h\in H_0$. Inclusion of $G_0$ into $HG_0$
thus induces a bijection of coset sets
$G_0/(H_0U)\simeq HG_0/HU$. The index formula along
$HU\subseteq HG_0\subseteq G$ consequently gives
\[
 a_H(u)=[G:HG^u]=f[G_0:H_0G^u]=f b_H(u).
\]
It follows that $n-a_H(u)=f(e-b_H(u))$ for every $u\geq0$.
Multiply the preceding integral identity by $f$ and use $n=ef$ to
obtain the general formula. In the totally ramified case $f=1$,
which gives its stated specialization. This argument combines the
cited different identities with the subgroup Herbrand integral;
it retains $|H_0|=e(L/M)$ rather than replacing it by $|H|$.
\end{proof}

\begin{lemma}[Finite upper orbits and the intrinsic APF index]
\label{fw:finite-upper-orbit-apf-index}
\uses{fw:finite-nonnormal-apf,fw:subextension-different-orbit-integral}
\lean{FontaineWintenberger.finite_upperOrbitIndex_eq_zeroIndex_mul_apfIndex}
\leanok
For every $u\geq0$, the actual finite upper orbit index satisfies
\[
 [G:HG^u]=[G:HG_0]\,a_{M/K}(u),
\]
where $a_{M/K}$ is the intrinsic APF inertia index. In particular,
if $M/K$ is totally ramified, the two indices agree.
\end{lemma}
\begin{proof}
Restriction of the absolute upper group to the finite Galois
extension identifies the intrinsic APF inertia index with
$[G_0:H_0G^u]$. The upper-orbit factorization proved above gives
$[G:HG^u]=[G:HG_0][G_0:H_0G^u]$, proving the formula.
The zero orbit index $[G:HG_0]$ is the actual residue degree
$f(M/K)$, which is one for a totally ramified extension.
\end{proof}

\begin{corollary}[The index profile of a prime-degree subextension]
\label{fw:prime-degree-nonnormal-index-profile}
\uses{fw:subextension-different-orbit-integral,fw:finite-upper-orbit-apf-index,fw:finite-first-upper-break,fw:herbrand-functions}
\lean{FontaineWintenberger.exists_prime_index_upperOrbitIndex_profile, FontaineWintenberger.subextension_prime_degree_different_profile, FontaineWintenberger.intermediate_spectral_prime_degree_different_profile, FontaineWintenberger.apfIndex_prime_degree_different_profile_of_finiteGaloisIntermediateField}
\leanok
Suppose $M/K$ is totally ramified of prime degree $p$, still without
assuming it normal. There is a finite real number $b\geq0$ such that
\[
 a_H(u)=
 \begin{cases}
 1,&0\leq u\leq b,\\
 p,&u>b,
 \end{cases}
 \qquad
 d(M/K)=(p-1)(b+1).
\]
The same profile holds for the intrinsic APF index of $M/K$.
\end{corollary}
\begin{proof}
For every normal subgroup $U$ of $G$, the index $[G:HU]$ divides
$[G:H]=p$, by the index formula for $H\subseteq HU\subseteq G$.
It is therefore either $1$ or $p$. At $u=0$, the index formula gives
$[HG_0:H]=[G_0:H_0]=e=p=[G:H]$, so $HG_0=G$ and $a_H(0)=1$.
The upper groups decrease, hence
$a_H$ is nondecreasing; eventual triviality makes its eventual value
$p$. On every bounded interval there are only finitely many upper
breaks, namely images of integer lower breaks under $\varphi$.
The convention $G_t=G_{\lceil t\rceil}$ and continuity of
$\varphi^{-1}$ include the right endpoint of each interval of
constancy. Hence the set where $a_H=1$ is a bounded initial interval
with its endpoint included. Write it $[0,b]$; this also covers $b=0$.
The orbit integral is now $(p-1)b$. The preceding theorem gives
$d(M/K)=p-1+(p-1)b$.

The finite nonnormal APF index is the index of the product of the
fixing subgroup with the actual upper group, relative to the same
product at zero. The latter product is $HG_0=G$ here. Its index is
therefore exactly $[G:HG^u]=a_H(u)$, which proves the last assertion.
\end{proof}


\begin{lemma}[Invariance of the actual ramification index and different order]
\label{fw:finite-extension-e-d-invariance}
\uses{fw:ramification-field-equiv,fw:spectral-valuation-ring-unique,fw:chosen-extension-integers,fw:integral-generator,fw:monogenic-different,fw:ideal-order-generator}
\lean{FontaineWintenberger.spectral_ramificationIdx_algEquiv, FontaineWintenberger.spectral_different_order_algEquiv, FontaineWintenberger.extensionValuationSubring_ramificationIdx_eq_spectral, FontaineWintenberger.extensionValuationSubring_different_order_eq_spectral}
\leanok
Let $L/K$ be finite separable. The ramification index and the
normalized order of the different of its actual integer ring are
preserved by every $K$-algebra isomorphism of $L$ with another
finite extension. They are also unchanged if its spectral
valuation is replaced by any compatible extension valuation.
\end{lemma}
\begin{proof}
Uniqueness of the integral closure identifies the two actual
integer rings by an isomorphism over $\mathcal O_K$; for a field
isomorphism this is its actual restriction. This isomorphism
preserves the maximal ideal and its powers, hence normalized
orders of elements. If $\pi$ is a base uniformizer, its order in
the extension ring is the ramification index; preservation of
this order proves the first assertion. Choose an integral
generator $\theta$ of the original extension ring. Its image
still generates, because the ring isomorphism is surjective.
The minimal polynomials over $\mathcal O_K$ agree, and derivative
evaluation commutes with the isomorphism. The monogenic different
formula expresses each different as the principal ideal of the
corresponding derivative value. Its normalized order is the order
of that value, which is preserved. This proves both invariance
claims without comparing the raw value groups of the valuations.
\end{proof}

\begin{theorem}[The intrinsic profile of an abstract prime-degree extension]
\label{fw:intrinsic-prime-degree-apf-profile}
\uses{fw:prime-degree-nonnormal-index-profile,fw:finite-extension-e-d-invariance,fw:finite-galois-overfield-apf,fw:apf-index-embedding,fw:apf-herbrand-homeomorphism}
\lean{FontaineWintenberger.intermediate_apfIndex_prime_degree_different_profile, FontaineWintenberger.apfIndex_prime_degree_different_profile, FontaineWintenberger.apfPsi_eq_of_index_profile, FontaineWintenberger.apfPhi_eq_of_index_profile, FontaineWintenberger.apf_prime_degree_different_profile, FontaineWintenberger.apfIndex_prime_degree_eq_different}
\leanok
Let $L/K$ be any finite separable totally ramified extension of
prime degree $\ell$, with no normality assumption. If $d$ is its
actual different order, put $b=d/(\ell-1)-1$. Then $b\geq0$ and
\[
 a_{L/K}(u)=\begin{cases}1&0\leq u\leq b,\\ \ell&u>b,\end{cases}
 \qquad
 \psi_{L/K}(u)=\begin{cases}u&u\leq b,\\b+\ell(u-b)&u>b,\end{cases}
\]
for every real argument in the indicated ranges. Its inverse is
\[
 \varphi_{L/K}(v)=\begin{cases}v&v\leq b,\\
 b+(v-b)/\ell&v>b.\end{cases}
\]
In particular $d=(\ell-1)(b+1)$; the negative continuations of
both functions are the identity.
\end{theorem}
\begin{proof}
Embed the original separable field into a separable closure and
take its finite normal closure $\Omega/K$. Its image $M$ is
isomorphic to $L$ over $K$, has degree $\ell$, and its actual
ramification index and different order agree with those of $L$
by the preceding invariance lemma. Apply the prime-degree orbit
profile to $M\subseteq\Omega$. Total ramification gives residue
degree one, so the inertia-relative APF index equals the full
orbit index. Invariance of the intrinsic APF index under the
original field isomorphism gives a break $c\geq0$ for $L$ and
$d=(\ell-1)(c+1)$. Since $\ell-1>0$, this identity gives $c=b$.

For $0\leq u\leq b$, integrating the index $1$ from zero gives
$\psi(u)=u$. For $u>b$, split the integral at $b$; the value at
$b$ is $b$ and the remaining constant index is $\ell$, giving
$b+\ell(u-b)$. For $u\leq0$ the prescribed continuation is the
identity, so the same piecewise formula holds on all real
arguments. The displayed formula for $\varphi$ follows by
substitution in these two expressions for the strictly increasing
function $\psi$ and uniqueness of its inverse. This also checks
the value at the closed breakpoint and at zero.
\end{proof}
